\section{Related work}
\label{sec:related}

We review computational work on Canadian Fish and related methods for hidden-state inference,
imperfect-information search, team coordination, and agent evaluation. For each group, we
distinguish methods used in this study from alternatives requiring assumptions or computational
resources not adopted here.

\subsection{Prior work on Canadian Fish}
\label{sec:related-fish}

The published material on this game is human and informal, and neither of its two substantial
sources is quantitative or describes a program. McLeod's Pagat entry \cite{pagat} is the canonical
rules source and records several points on which published variants differ; ours is specified in
\S\ref{sec:game-dialect}. Develin's chapter \cite{develin} is the most detailed strategy writing we
located, describes the same variant, and supplies two observations this design uses: an ask teaches
the opponents as much as it teaches a teammate, and a player is better off not declaring a half-suit
it holds entirely, since holding one carries no risk. The second is a tactic and not a rule: the
engine imposes no such condition, and \S\ref{sec:game-rules} states what it does impose. Secondary strategy pages \cite{dorsa,strategy-site,wikiliterature} restate blackballing,
asking the asker, and an endgame delegation heuristic; they are evidence about how the game is
played and talked about, not sources of quantitative claims.

\paragraph{Computational prior work.} Literature searches conducted during this project's
development cycles, covering academic indexes, public code repositories and the RLCard toolkit
\cite{rlcard}, found no academic work on the six-player game and no public repository containing
search, planning or a working learner for it. The one agent located, Somani's \texttt{literature}
\cite{somani,somani-server}, plays the four-player, 48-card variant, which does not raise the
six-player allocation problem, and publishes no strength numbers; a closed-source mobile
application advertises bots with no stated methodology \cite{playstore}. Every novelty statement in
this paper therefore takes the form ``we found no prior published work that \ldots{}'' and is a
statement about those searches rather than a demonstrated absence. The Valet testbed
\cite{goadrich2026valet} packages twenty-one traditional imperfect-information card games in a
common description language with measured branching factors and durations; a shared testbed of
that kind would supply the external calibration this line of work lacks, but no comparison against
it is attempted here.

\subsection{Determinization and perfect-information analysis}
\label{sec:related-determinization}

Determinization, or perfect-information Monte Carlo (sample hidden cards consistently with the
public history, solve the resulting perfect-information game, and play the move with the best average
value), is the standard recipe for hidden-hand card games \cite{levy,ginsberg-gib}, and is
unsound however many worlds are drawn. Frank and Basin named the two mechanisms: \emph{strategy
fusion}, in which the searcher plays a different strategy in each sampled world although a real
player must commit to one strategy per information set, and \emph{non-locality}, in which a node's
value depends on parts of the tree outside its own subtree
\cite{frank-basin-ai,frank-basin-aaai}. Long, Sturtevant, Buro and Furtak characterise when the
method nevertheless performs well, through leaf correlation, bias and a disambiguation factor
\cite{long-pimc}; Information Set Monte Carlo Tree Search replaces independent trees with a single
tree over information sets \cite{cowling-ismcts,goodman-ris}; the $\alpha\mu$ family backs up Pareto
fronts of world-vectors instead of scalar averages, repairing both mechanisms, and reports Bridge
3NT declarer play at 60.2\% to 63.0\% with twenty worlds \cite{cazenave2019alphamu}; and
knowledge-based paranoia search reasons about the worst case over deduced-consistent worlds rather
than the average, reporting over 1{,}000 points on the extended Seeger scale in Skat
\cite{edelkamp2021paranoia}.

The evaluated variant presents an additional limitation for clairvoyant rollout evaluation, stated
in \S\ref{sec:game-why}: under clairvoyant play the available asks are scarcely distinguished in
value, so repairing strategy fusion and non-locality, as $\alpha\mu$ does, would not by itself make
a clairvoyant rollout informative here. We have not tested $\alpha\mu$ or its relatives on this
game, and the statement is about the evaluated variant rather than about those methods.

What this corpus adds is that the binding problem is not the abstract failure but \emph{the
statistic}. One determinization's value is the final half-suit differential, whose spread is an
order of magnitude larger than the differences between the leading candidates, so an argmax over
sampled means selects on residual noise, and does so more confidently the more candidates it is
offered. That is the optimiser's curse rather than strategy fusion, and \S\ref{sec:background-vsix}
reports it measured at \vsixCurseGap{}~pp by a control that forces the same code to return the
blueprint's own pick. These experiments motivate the guarded action-selection rule in
\S\ref{sec:agent-search}.

\subsection{Search under imperfect information}
\label{sec:related-search}

States in an imperfect-information game do not have well-defined values, so depth-limited search
needs something at the leaf that represents the opponent's ability to choose. Brown, Sandholm and
Amos supply it: let the opponent select among $k$ continuation strategies at the depth limit, each
yielding its own leaf-value vector, and require the searcher to be good against that choice
\cite{brown2018depthlimited}. The soundness argument is stated for two-player zero-sum games, and
Fish's team-level game presents three uncorrelated opponent seats at the leaf, so nothing is
inherited; it is adopted here as a robustification heuristic with a measurable consequence, and
audited rather than assumed.

Two lines address what a search may prune. Zhang and Sandholm solve subgames over low-order
knowledge rather than the full common-knowledge closure, and prove that the naive form can
\emph{increase} exploitability \cite{zhang2021subgamesolvingwithoutck}. That warning applies directly
to the live-ask gate this policy family ships, which is a knowledge-limited pruning rule, and it is
why the side-channel and pathology gates of \S\ref{sec:evaluation-gate} were built before any
search work rather than after it. Subgame solving in adversarial team games builds a gadget over the
team belief DAG and runs in polynomial time given a best-response oracle under a sparsity condition
on the blueprint \cite{zhang2022teamsubgame}; the cited sparsity condition is not established for
the blueprint used here, and we did not attempt the construction.

In cooperative games the constraint transfers verbatim. SPARTA improves a blueprint purely at test
time and distinguishes single-agent search, which assumes every other agent strictly follows the
blueprint, from multi-agent search, in which all agents run the \emph{same} common-knowledge
procedure \cite{lerer2020sparta}. The moment a second teammate also searches, the first teammate's
belief over its partner is computed under a policy that partner is no longer playing, and the
improvement guarantee is void. Fish satisfies the requirement cheaply, for the reason
\S\ref{sec:game-why} gives; the qualifier about randomisation is what
\S\ref{sec:background-vsix} turns into an engine switch.

The full-solver line \cite{moravcik2017deepstack,brown2020rebel,schmid2023sog} established that
test-time re-solving over public belief states is supported by the formal models and assumptions in
those studies. We do not implement these algorithms, because each requires iterating or learning
over a closure whose cardinality at deal time is $45!/(9!)^5 \approx 1.9 \times 10^{28}$ deals from
one seat. Update equivalence recasts search as replicating a last-iterate-convergent update locally,
needing no public-information enumeration, and reports substantially lower search time than
public-belief-state search in Hanabi \cite{sokota2024updateequivalence}; its improvement guarantee
is stated for common-payoff games, which Fish is not. Learned Belief Search \cite{hu2021lbs} learns
a belief model for settings in which the exact belief is unavailable; we do not use it because the
constraint-conditioned belief is computed exactly here and the behavioural weighting is deliberately
simple, not because a learned belief could not add information. Online planning by policy
fine-tuning \cite{fickinger2021rlfinetuning} requires gradient updates at decision time and was not
attempted in this CPU-only engine with no network. The depth-limited counter-strategy line
\cite{milec2021cdr,milec2025abd} carries one warning this study honours, robust-response methods
do not compose with limited-look-ahead solving off the shelf, and the learned-model and
gadget-refinement work \cite{kubicek2025lamir,kubicek2026refinements} is recorded for its
meta-lesson rather than its machinery: when a search subproblem has many near-equivalent optima,
which one is selected is not a tie-break. That is the same observation \S\ref{sec:background-vsix}
makes about an argmax resolved by array order.

The result that a search of this kind can \emph{increase} exploitability is Zhang and Sandholm's
\cite{zhang2021subgamesolvingwithoutck}. Where the other work above carries a guarantee at all, its
conditions and this game's fit to them are recorded above, method by method. What none of them
supports is reading a win rate as a bound on exploitability, which is why this paper measures the two
separately. Local best response gives an approximate lower bound on best-response value: a small greedy
search looking one action ahead and assuming checks thereafter showed each of the 2016 competition
entrants it evaluated to be more than 3{,}180\,mBB/h exploitable \cite{lisy2017lbr}, the
competition's winning submission was not among them, and the authors say why. Its two documented
pitfalls are carried into \S\ref{sec:results-adversarial}: restricting where the response may begin
deviating understates exploitability, and a negative score proves nothing. Accordingly every
adversarial-response figure in this paper is reported as the observed win-rate edge of a particular
response within its class, which is a lower bound on the best-response value over that class only up
to sampling error, and never as an upper bound.

\subsection{Policy-aware inference}
\label{sec:related-inference}

Skat engines enumerate the deals consistent with the public record and reweight them by a model of
what each player would have done, a supervised per-card location model
\cite{solinas-supervised}, or reach probabilities taken from a policy \cite{rebstock-inference}, and a dedicated inference module was reported worth $+217$ tournament points per 36 games to the
engine Kermit on its own \cite{buro-kermit}. The policy prior of \S\ref{sec:inference-prior} is the
same idea in a much simpler form: a fitted soft reweighting of the rules posterior, with two
parameters and no learned opponent model.

This is the literature that has measured the same trade-off, and it is where
\S\ref{sec:background-vsix}'s belief-model comparison belongs. The finding there is that in this
game the constraint-conditioned distribution is computable exactly by a tractable factorised
dynamic program and is cheap, so in this implementation computational cost is not the only
motivation for approximation, and yet substituting the exact constraint-conditioned distribution
into an otherwise identical policy loses to the approximation with behavioural weighting, on two
independent seed banks.

Two adjacent lines are named. Belief modelling in cooperative games, the Bayesian Action Decoder's
public-belief formulation \cite{foerster2019bad} and its factorised-prescription successor
\cite{sokota2021capi}, is aimed at games whose exact belief is intractable, which is not the
situation here. The permanent-like counts that a constraint system of this shape generally induces
are the subject of exact formulae, Ryser's inclusion--exclusion formula \cite{ryser} and Glynn's
sign-vector formula \cite{glynn}, and of approximation algorithms, the Jerrum--Sinclair--Vigoda
chain \cite{jsv}, matrix scaling \cite{lsw} and sequential importance sampling \cite{chen-sis}; the
block structure of \S\ref{sec:inference} makes the exact count tractable directly, so none of the
five is needed for it. Matrix scaling appears in the deployed path for a different reason: the
approximate belief replaces the exact count join with a Sinkhorn iteration of that kind, not because
the exact count is unaffordable but to incorporate the two behavioural weighting parameters
(\S\ref{sec:inference-prior}). No learned belief model appears anywhere in this corpus; that is a
stated non-goal of the project.

\subsection{Team games and common information}
\label{sec:related-team}

The formal statement of three teammates sharing a payoff, holding different information and unable
to communicate is the common-information approach \cite{nayyar2013common}, discussed in
\S\ref{sec:game-why}. The team-game equilibrium literature is cited to delimit scope: the target
concept for a Fish team is a team-maxmin equilibrium with correlation \cite{celli-gatti}, connected
to extensive-form correlation by Farina and coauthors \cite{farina-collusion} and made computable
through the team belief DAG and the two-player conversion \cite{zhang-tbdag,carminati-tpi}, with
subsequent work on hidden-role settings \cite{carminati2024hiddenrole}. None of it is computed here.
We have not established whether a Fish team satisfies the recall conditions those constructions
require; the evident obstacle is that each seat privately observes nine cards its teammates do not.

The cooperative-play line is cited for its evaluation discipline rather than its algorithms. Hanabi
was proposed as a benchmark precisely because self-play scores overstate coordination
\cite{bard-hanabi}, and the ad-hoc-teamplay results that followed
\cite{hu-sad,hu-otherplay,hu2021obl,hu2022sed,jacob2022pikl,hu2022humanregularized,dizdarevic2025ah2ac2}
are the reason the cross-play measurement of \S\ref{sec:results-crossplay} exists at all: a
convention that collapses when partners are swapped is brittleness, and it is invisible to self-play
evaluation. Two neighbouring results are cited for what they show about porting rather than as
ad-hoc teamplay: joint policy search is the bridge-bidding analogue, scoring simultaneous changes at
several information sets \cite{tian2020jps}, and the Hanabi construction that most clearly fails to
port is the hat-guessing code \cite{cox-hanabi}, which works because each receiver sees every other
receiver's hand, and in Fish no player sees any other player's hand. The same lesson from a different direction is Cicero's use of a human-regularised policy
\cite{fair2022cicero}.

\subsection{Evaluation methodology}
\label{sec:related-eval}

The evaluation framework of \S\ref{sec:evaluation} draws on duplicate-match and
variance-reduction methods. Every cell in this paper is a paired duplicate rather than an
independent sample, on the model of duplicate bridge and of the Annual Computer Poker Competition's
duplicate matches and common seeds \cite{acpc}, whose protocols are also the model for a fixed
opponent panel scored identically for every arm. The estimator-side alternatives, DIVAT \cite{zinkevich-divat} and AIVAT
\cite{burch-aivat}, subtract a learned or hand-built value function from a single sample instead;
neither is used here, and the pairing is what makes the design competitive with them without
introducing a fitted object. Abstraction pathology
\cite{waugh-abstraction} is the standing warning that a model can be improved on its own metric and
made worse in play; \S\ref{sec:development-leaf} illustrates a related mismatch between a surrogate
metric and playing performance. Re-evaluation of agent
populations \cite{balduzzi-reeval} and the rating literature
\cite{coulom-whr,herbrich-trueskill,nelo,nelo-fishtest} inform why this paper reports a worst case and a minimax
regret over a shared panel rather than a single rating number. Recent work on gold-standard
evaluation practice \cite{kelidari2026goldstandard} is consistent with the preregistration and
sealing discipline of \S\ref{sec:evaluation-seal}, which this project adopted independently.
